% =====================================================================
% GLOSSÁRIO (elemento pós-textual, depois das Referências - ABNT NBR 14724)
% Para adicionar um termo, copie uma linha \item e mantenha a ordem alfabética.
% Siglas (API, HTTP, JSON, REST etc.) ficam na Lista de abreviaturas.
% =====================================================================
\chapter*{Glossário}
\addcontentsline{toc}{chapter}{Glossário}
\markboth{Glossário}{Glossário}

\begin{description}
  \item[Angular] Framework de código aberto mantido pelo Google para a criação de aplicações web em TypeScript, organizadas em componentes.
  \item[Arduino] Plataforma de hardware e software de código aberto. O seu framework e a sua IDE também são usados para programar outras placas, como o ESP32.
  \item[Backend] Parte do sistema executada no servidor, responsável pelas regras de negócio, pelo acesso aos dados e pela comunicação com os clientes.
  \item[Biblioteca] Conjunto de código pronto que pode ser reaproveitado por outros programas.
  \item[Cabeçalho (header)] Campo de uma requisição ou resposta HTTP que leva informações adicionais, como credenciais de acesso e o tipo do conteúdo.
  \item[Carregamento sob demanda (lazy loading)] Técnica em que partes de uma aplicação só são carregadas quando o usuário precisa delas.
  \item[Chart.js] Biblioteca JavaScript de código aberto para a criação de gráficos em páginas web.
  \item[Cloud Firestore] Banco de dados NoSQL em nuvem do Firebase, que armazena os dados em documentos organizados em coleções.
  \item[Código aberto] Modelo em que o código-fonte de um software é disponibilizado publicamente e pode ser estudado, modificado e redistribuído conforme a sua licença.
  \item[Componente standalone] No Angular, componente que declara as próprias dependências, sem depender de um módulo (NgModule).
  \item[Controlador (controller)] Classe da API que recebe as requisições HTTP de uma rota, repassa o trabalho aos serviços e devolve a resposta.
  \item[CORS] Cross-Origin Resource Sharing. Mecanismo dos navegadores que controla se uma página carregada de um domínio pode fazer requisições a um servidor de outro domínio; o servidor indica, por cabeçalhos, quais origens são permitidas.
  \item[Docker] Plataforma que empacota uma aplicação e as suas dependências em uma imagem, executada em contêineres isolados, de forma igual em qualquer máquina.
  \item[DTO] Data Transfer Object. Objeto que define o formato dos dados trocados entre o cliente e a API, separado do modelo armazenado no banco.
  \item[Endpoint (rota)] Endereço da API que, combinado a um método HTTP, executa uma operação específica.
  \item[Firebase] Plataforma do Google que oferece serviços prontos para aplicações, como banco de dados, autenticação e hospedagem.
  \item[Firebase Authentication] Serviço do Firebase que gerencia o cadastro e o login de usuários e emite os tokens de identificação.
  \item[Firmware] Programa gravado na memória de um dispositivo embarcado, que controla diretamente o seu hardware.
  \item[Framework] Conjunto de bibliotecas e convenções que fornece a estrutura básica de uma aplicação, sobre a qual o desenvolvedor escreve o próprio código.
  \item[FreeRTOS] Sistema operacional de tempo real de código aberto para microcontroladores, que permite executar várias tarefas de forma concorrente.
  \item[Frontend] Parte do sistema com a qual o usuário interage diretamente; no caso de uma aplicação web, é executada no navegador.
  \item[Git] Sistema de controle de versão distribuído, que registra o histórico de alterações do código-fonte.
  \item[GitHub] Plataforma on-line de hospedagem de repositórios Git.
  \item[GPIO] General Purpose Input/Output. Pino de uso geral de um microcontrolador, que pode ser configurado como entrada ou saída.
  \item[Guard] No Angular, função executada antes de abrir uma rota da aplicação, que decide se o usuário pode acessá-la.
  \item[Interceptador (interceptor)] Função que atua sobre todas as requisições HTTP de uma aplicação, antes do envio ou após a resposta, por exemplo para incluir um token ou tratar erros.
  \item[Java] Linguagem de programação orientada a objetos, executada sobre uma máquina virtual (JVM).
  \item[JUnit] Framework de testes unitários para Java.
  \item[Long-polling] Técnica em que o cliente faz uma requisição e o servidor só responde quando há um dado novo ou quando se esgota um tempo limite; em seguida, o cliente abre uma nova requisição.
  \item[Maven] Ferramenta de compilação e de gerenciamento de dependências de projetos Java.
  \item[Mock] Objeto ou dado que simula o comportamento de uma dependência real durante os testes ou o desenvolvimento.
  \item[Mockito] Biblioteca Java para a criação de mocks em testes.
  \item[Mutex] Mecanismo de exclusão mútua que garante que apenas uma tarefa por vez acesse um recurso compartilhado.
  \item[NoSQL] Categoria de bancos de dados que não usam o modelo relacional de tabelas, armazenando os dados em documentos, pares chave-valor, colunas ou grafos.
  \item[OpenAPI] Especificação padrão para descrever as rotas, os parâmetros e as respostas de uma API REST em um arquivo legível por máquinas.
  \item[Render] Plataforma de nuvem para hospedagem e execução de aplicações web e APIs.
  \item[Repositório] No Git, pasta de um projeto com todo o seu histórico de versões. Na arquitetura da API, classe responsável por ler e gravar os dados no banco, isolando o resto da aplicação dos detalhes de acesso.
  \item[SDK] Software Development Kit. Conjunto de bibliotecas e ferramentas fornecido por um fabricante para integrar uma aplicação aos seus serviços.
  \item[Serviço (service)] Classe da API que concentra as regras de negócio e coordena os repositórios.
  \item[Signal] No Angular, valor reativo que atualiza automaticamente as partes da tela que dependem dele quando é alterado.
  \item[Spring Boot] Framework Java para a criação de aplicações e APIs web, com configuração automática e servidor embutido.
  \item[Spring Security] Módulo do Spring responsável pela autenticação e pela autorização das requisições.
  \item[springdoc-openapi] Biblioteca que gera automaticamente a especificação OpenAPI e a interface Swagger de uma aplicação Spring Boot.
  \item[Swagger] Conjunto de ferramentas que gera, a partir da especificação OpenAPI, uma página interativa com a documentação das rotas de uma API, na qual elas podem ser testadas pelo navegador.
  \item[Tailwind CSS] Framework CSS baseado em classes utilitárias aplicadas diretamente nos elementos HTML.
  \item[Thread] Fluxo de execução independente dentro de um programa, que pode ser executado em paralelo a outros.
  \item[Token] Sequência de caracteres emitida após a autenticação, que comprova a identidade do usuário e é enviada nas requisições seguintes.
  \item[TypeScript] Linguagem que estende o JavaScript com tipagem estática e é convertida em JavaScript para ser executada.
  \item[Variável de ambiente] Valor configurado no ambiente em que a aplicação é executada, fora do código-fonte, usado para guardar dados como chaves de acesso e endereços.
  \item[Webhook] Rota de um servidor que recebe dados enviados automaticamente por outro sistema quando ocorre um evento. No sistema, são as rotas pelas quais o ESP32 envia as leituras, o estado e o histórico.
\end{description}
